% Part: methods
% Chapter: induction
% Section: introduction

\documentclass[../../../include/open-logic-section]{subfiles}

\begin{document}

\olfileid{mth}{ind}{int}

\olsection{પરિચય}

આગમન સાબિતીની મહત્વની પદ્ધતિ છે. તે જુદાં જુદાં સ્વરૂપે
તર્કશાસ્ત્ર, સૈદ્ધાંતિક કમ્પ્યુટર વિજ્ઞાન અને ગણિતના લગભગ દરેક
ક્ષેત્રમાં વપરાય છે. તર્કશાસ્ત્રનાં ઘણાં પરિણામો સાબિત કરવા તેની
જરૂર પડે છે.

આગમનને ઘણી વાર નિગમનથી જુદું ગણવામાં આવે છે અને વિશેષ કિસ્સાઓમાંથી
સામાન્ય દાવા સુધી પહોંચતા અનુમાન તરીકે વર્ણવાય છે. દાખલા તરીકે,
આપણે ઘણા લીલા મરકત જોયા હોય અને જેને મરકત કહીએ એવું એક પણ
બિનલીલું ન જોયું હોય, તો બધા મરકત લીલા છે એવું તારવી શકીએ. આ
આગમનાત્મક તારણ છે: તે ઘણા વિશેષ કિસ્સાઓમાંથી (આ મરકત લીલું છે,
પેલું મરકત લીલું છે વગેરે) સામાન્ય દાવા (બધા મરકત લીલા છે) સુધી
પહોંચે છે. \emph{ગાણિતિક} આગમન પણ સામાન્ય દાવો સ્થાપે છે, પરંતુ
તે આ "સરળ આગમન"થી ઘણું જુદા પ્રકારનું છે.

ખૂબ સાદા શબ્દોમાં, ગણિતમાં આગમનથી કરેલી સાબિતી બતાવે છે કે
ચોક્કસ પ્રકારના બધા ગાણિતિક પદાર્થોમાં કોઈ એક ગુણધર્મ છે. સૌથી
સરળ કિસ્સામાં આ પદાર્થો પ્રાકૃતિક સંખ્યાઓ હોય છે. ત્યારે
આગમનથી કરેલી સાબિતી બધી પ્રાકૃતિક સંખ્યાઓમાં તે ગુણધર્મ છે
એવું સ્થાપે છે. તે આ રીતે કરે છે:
\begin{enumerate}
    \item $0$માં તે ગુણધર્મ છે, અને
    \item જ્યારે પણ કોઈ સંખ્યા~$k$માં તે ગુણધર્મ હોય, ત્યારે~$k+1$માં પણ હોય છે.
\end{enumerate}
પ્રાકૃતિક સંખ્યાઓ પરનું આગમન ઘણી વાર એવી ગાણિતિક વસ્તુઓ અંગેના
સામાન્ય દાવા સાબિત કરવા પણ વાપરી શકાય છે જેમને સંખ્યા આપી શકાય.
દાખલા તરીકે, દરેક સીમિત ગણમાં ઘટકોની સીમિત સંખ્યા~$n$ હોય છે.
જો આગમનથી બતાવી શકીએ કે દરેક સંખ્યા~$n$ માટે "કદ~$n$ના બધા
સીમિત ગણો \dots છે", તો બધા સીમિત ગણો વિશે કંઈક સાબિત થયું હશે.

આગમનને \emph{આગમનાત્મક રીતે વ્યાખ્યાયિત} ગાણિતિક પદાર્થો સુધી
પણ વિસ્તારી શકાય છે. દાખલા તરીકે, પ્રથમ-ક્રમ તર્ક જેવી
ઔપચારિક ભાષાની અભિવ્યક્તિઓ આગમનાત્મક રીતે વ્યાખ્યાયિત
થાય છે. \emph{સંરચનાત્મક આગમન} એવી બધી અભિવ્યક્તિઓ અંગે
પરિણામો સાબિત કરવાની રીત છે. ખાસ કરીને સંરચનાત્મક આગમન
તર્કશાસ્ત્રમાં ખૂબ ઉપયોગી છે અને વ્યાપકપણે વપરાય છે.

\end{document}
